\chapter{ANALYSIS OF PERFORMANCE TESTING RESULTS}
\label{chap:analysis_of_results}
Diskusjon rund resultatene. Det skal gjøres en evaluering av hvorvidtvalgt løsning. Avvik drøftes.
det er samsvar mellom kravspesifikasjonen satt opp og

An update package usually contains only a small number of changes from the
previous version

\section{Effectiveness of update processing}
It came as a big surprise to us that \bhtb{} performed so well in the
update tests after our client had listed that as one of it's major faults. We
suspect that it is the very simple nature of it's data structure that is making
the difference. We also suspect that our tests didn't really reflect on the
reality of our clients setup and database complexity, as they have multiple
sensors and many rules turned on or off, something our tests didn't simulate. 

\snortmanager{} also did well in the tests, which wasn't really all that
surprising, as it has a very similar data structure to \bhtb{} in that it is
very simple. But it seems to start having more problems once we introduced more
files into the tests. There is a big difference between tests \textit{SYN-1000}
and \textit{SYN-1000-B}, where we have the same number of rules, but spread over
more files. This problem does not seem to exist in either \snowman{} nor
\bhtb{}.

None the less, both \snortmanager{} and \bhtb{} pummelled \snowman{} v0.4 into
the ground in these tests, which led us to completely rework our update program
workflow. Our problem was that for every rule, we were writing its data to the
database and then moving on to the next rule. This caused a lot of busy-waiting
for the database. You can see that the times accelerated a lot on the machine
with the lower end harddrives, so the problem also manifests itself more when
the harddrives can't keep up.

We proceeded to rewrite the code so that all the rules are parsed into
memory first and then is committed into the database in one big query. As you
can see in table \ref{tabell med tall} we improved \snowman{}'s performance
in v0.5 with as much as xx\% in some tests. 

While we still didn't fulfil the performance goal, considering the difference in
complexity and structure, we feel the difference is acceptable. The results in
this test is a very good indicator that performance testing during development
is a very useful thing to a development project.

\section{Effectiveness of rule distribution}
As we expected, \bhtb{} had a very flat performance curve in this test, but not
exactly flat. There is a small increase in time spent when the number of rules
go up, as you can see in table \ref{tabell}, meaning there is some degree of
processing time infuencing the results. 

Still, considering the tests on S1 and
S2 were carried out on the same network and sensor, and S3 were carried out on a
completely different network and sensor, it is apparent that distribution from
\bhtb{} is mostly determined by how fast the data can be moved from the server
to the sensor. 

The same difference in testing environments is apparent in the \snowman{} tests.
S1 and S2 have pretty much the same results while S3 probably benefits from it's
different hardware and network setup. This leads us to conclude that the
performance of distribution is dictated by the bandwidth and by the hardware of
the sensor, and not so much the hardware of the central server.

However, there does not seem to be a pure performance gain over \bhtb{} as we
expected by going for a distributed model with a change based distribution
protocol. But we are seeing a significant reduction in time spent distributing
when there are only updates and not a massive batch of rules. While \bhtb{} uses
roughly the same time on tests SF-28064-A and SF-28064-B, we use about 80\% less
time on test B than on test A. 

This gives merit to the model we picked for our distribution architecture and
with a little more performance optimizations, the overall times might improve
further as well.

\section{GUI latency}
These tests probably entails the most grim part of our performance results as we
seeminly have failed to produce an improvement over \bhtb{} when it comes to
displaying rules, something that does not make much sense if you do the math. We
are only fetching and displaying a fraction of the number of rules that \bhtb{}
does\footnote{About 100 to 28000.}, yet we spend about twice as long on it. 

We have not spent much time investigating the cause of this, but we strongly
suspect its a combination of our more complex data structure, coupled with our
rather novice use of a complex template framework such as Django.

A small optimization was attempted by enabling a setting that lets Django
prefetch related objects from the database and that did improve the results
somewhat, but not enough to be noticable. 

\section{Linearity in performance}
Må ha grafer for å snakke om dette.

\section{Conclusions}
- Snowman kom dårligere ut
- Fordi større data struktur
- Optimaliseringen i update var god
- Viser viktigheten av performance testing i prosjekter

